% Liam Solus, Polyhedral Problems in Combinatorial Convex Geometry,
% University of Kentucky (2015), DOI 10.13023/ETD.2016.014.
% https://uknowledge.uky.edu/math_etds/32/
% Definition 2.1.1, p. 23; verbatim conjecture on printed p. 73.
% Also Braun and Solus, r-Stable hypersimplices, JCTA 157 (2018),
% 349--388, Definitions 1.1--1.2 and Conjecture 4.22,
% https://doi.org/10.1016/j.jcta.2018.03.002 .
\documentclass{article}
\usepackage{amsmath,amssymb}
\newcommand{\rstablentwo}{\Delta_{n,2}^{stab(r)}}
\begin{document}
Label the regular $n$-gon cyclically by $[n]=\{1,\ldots,n\}$ and put
$\operatorname{cd}(i,j)=\min(|i-j|,n-|i-j|)$.
For $0<k<n$ and $1\le r\le\lfloor n/k\rfloor$, an $r$-stable $k$-subset
$I\subseteq[n]$ satisfies $\operatorname{cd}(i,j)\ge r$ for all distinct
$i,j\in I$. Its characteristic vector is $\epsilon_I=\sum_{i\in I} e_i$
in the standard basis of $\mathbb Z^n$.

\noindent\textbf{Definition 2.1.1 (in formulas).}
\[
\Delta_{n,k}^{stab(r)}=\operatorname{conv}
 \{\,\epsilon_I:\ |I|=k,\ I\text{ is }r\text{-stable}\,\}.
\]
In particular $k=2$ means $e_i+e_j$, not $\pm(e_i-e_j)$.

\noindent\textbf{Conjecture 3.4.5.}
The $h^*$-polynomials of the $r$-stable second hypersimplices
$\rstablentwo$ are unimodal.

For a lattice polytope $P$ of affine dimension $D$ the normalization is
\[
\sum_{t\ge0}|tP\cap\mathbb Z^n|x^t =
\frac{h^*(P;x)}{(1-x)^{D+1}}.
\]
Unimodality means that for some index $j$, the coefficients are weakly
increasing up to $j$ and weakly decreasing thereafter.
\end{document}
